\ifdefined\EnableTaggedPDF
  \DocumentMetadata{lang=en-AU,pdfstandard=ua-2,tagging=on}
\fi
\documentclass[10pt]{article}

\usepackage[T1]{fontenc}
\usepackage{lmodern}
\usepackage[margin=1in]{geometry}
\usepackage{microtype}
\usepackage{booktabs}
\usepackage{array}
\usepackage{enumitem}
\usepackage[hidelinks]{hyperref}

\hypersetup{
  pdftitle={Detecting Strategically Decomposed Malicious Code Changes with Subscription-Backed LLM Reviewers},
  pdfauthor={Tom Howard},
  pdflang={en-AU},
  pdfdisplaydoctitle=true
}

\newcommand{\code}[1]{\texttt{#1}}

\title{Detecting Strategically Decomposed Malicious Code Changes\\
with Subscription-Backed LLM Reviewers}
\author{Tom Howard\\
\href{mailto:tom@windyroad.com.au}{tom@windyroad.com.au}}
\date{19 July 2026}

\begin{document}
\maketitle

\begin{abstract}
Large language models are increasingly used to review code changes, yet most
evaluations present a vulnerable function, file, or repository as one unit. A
malicious change can instead be decomposed across submissions that appear
individually routine but compose into an unsafe final state. We describe a
prospectively specified study of whether this decomposition reduces detection
and whether detection differs between pull-request gates and trunk workflows
that treat main as untrusted until pipeline review completes. The study uses 40
matched malicious-benign pairs drawn from 40 structural templates in eight
abstract policy families. Every generated case module is synthetic, executes
only in memory, and has no external capability in the audited fixed corpus.
Atomic and three-submission variants have identical final trees and changed-line
totals. A factorial base design varies intent, decomposition, and workflow
framing across two confirmatory subscription-backed systems: Codex CLI and
Claude Code. A prospectively selected 16-template nested stratum adds
cumulative-history context and a second trial. The full two-system schedule
contains 2,816 review boundaries, including the 1,280-boundary base stratum used
for H1--H3. A separately reported, terms- and preparation-contingent Ollama Cloud stratum
would reuse only 640 base boundaries and cannot enter the two-system estimands.
The primary endpoint is \code{blocked\_by\_activation}; paired intent
discrimination is the H1 contrast over that endpoint.
Workflow framing and the decomposition-by-context interaction are prespecified
secondary hypotheses; decomposition by workflow and the cumulative-context main
effect are exploratory. This pre-results
manuscript reports the design, analysis plan,
deterministic baseline, safety controls, and limitations; it contains no LLM
benchmark outcomes.
\end{abstract}

\section{Introduction}

Automated code review is commonly evaluated as a classification problem: a
reviewer receives one code artifact and predicts whether it contains a defect.
Real delivery systems expose a different unit of work. Reviewers see a sequence
of submissions, and an unsafe capability may emerge only when several changes
compose. Each earlier change can be necessary for the final failure without
being unsafe in isolation.

This distinction matters for both pull-request and trunk-based delivery. A
pull-request gate asks whether a proposed change may enter the protected branch.
A trunk workflow can permit a commit onto main while still treating main as
untrusted until deterministic and model-based pipeline gates complete. These
workflows differ in policy, but empirical comparisons often also change the
evidence shown to the reviewer.

We study whether two confirmatory subscription-backed LLM review products
recognise a composition before its activation boundary. A conditional Ollama
Cloud replication is declared separately and cannot alter confirmatory results.
The benchmark represents harm as an abstract policy assertion: an unprivileged
synthetic actor receives a protected synthetic capability. It neither implements
an exploit nor interacts with a real target. The contribution is a reproducible,
prospectively specified test of sequence-level review with no usage-based spend.

\section{Related Work}

Yu et al. evaluated seven LLMs and multiple prompt forms for security code
review, comparing them with static analysis and examining factors associated
with detection quality~\cite{yu2024}. VulDetectBench separates vulnerability
identification, classification, and localisation~\cite{liu2024}. Zhou et al.
compare repository-level LLMs and static analysis tools and report a
recall--false-positive trade-off~\cite{zhou2024}. We retain a deterministic
conventional-lint baseline but do not claim that lint and semantic security
review are equivalent tasks.

Field evidence also indicates that context and interaction design affect
developer trust and usefulness. Adalsteinsson et al. report context switching,
missing context, false positives, and trust as practical concerns in an
industrial study of LLM-assisted review~\cite{adalsteinsson2025}. Basic and
Giaretta synthesise evidence on LLM code-security benefits, false findings,
prompting, and failure modes~\cite{basic2024}. Our study changes the unit of
analysis from one artifact to a sequence of submissions.

\section{Research Questions}

The primary endpoint is \code{blocked\_by\_activation}: at least one review at or
before the prespecified activation boundary returns \code{block}.

\begin{description}[leftmargin=3.4em,style=nextline]
  \item[RQ1] Can the review systems discriminate malicious sequences from
  matched benign sequences?
  \item[RQ2] Does decomposition across three submissions reduce timely
  detection relative to an atomic submission with the same final tree?
  \item[RQ3] Do estimates differ between pull-request and untrusted-main trunk
  framing?
  \item[RQ4] Does supplying cumulative submission history change the malicious
  split penalty, averaged equally over workflows, and how consistent are repeated reviews?
\end{description}

H1 predicts that the malicious block rate exceeds the benign block rate, equally
averaged over decomposition, workflow, and the two confirmatory products. H2
predicts that the malicious block rate is lower for split sequences than for
atomic submissions, equally averaged over workflow and the two products. H3 is
nondirectional: among malicious base local trial-1 cells, trunk and pull-request
framing differ after equal averaging over atomic and split conditions and both
products. H4 predicts that the malicious-only nested interaction
\((\mathrm{split}_{\rm cumulative}-\mathrm{atomic}_{\rm cumulative})-
(\mathrm{split}_{\rm local}-\mathrm{atomic}_{\rm local})\), equally averaged
over pull-request and trunk after both trials and products are averaged within
template, is positive. H1 and H2 are primary; H3 and H4 are prespecified
secondary hypotheses. The decomposition-by-workflow interaction and
cumulative-context main effect are exploratory, and no workflow-equivalence
claim is tested.

\section{Method}

\subsection{Safe benchmark}

The full deterministic generator contains 200 structural templates formed by
eight scenario families crossed with five data representations and five
control-flow shapes. The active subscription-bounded study selects template
indexes 1, 7, 13, 19, and 25 from each family. This balanced diagonal retains
every representation and flow shape once per family, giving 40 structural
templates. Only the first invented-identifier instance is retained.

Each template contains a malicious-benign pair. Every case has an atomic
submission and a three-submission sequence. Atomic and split forms have
identical final source trees, total changed-line counts, activation timestamps,
and policy outcomes. For malicious sequences, the generator reverts each split
step separately and executes the oracle; all three counterfactuals must remain
safe. Benign cases match structure and salience but do not grant the protected
capability.

Generated case modules import only relative synthetic modules and run in memory.
Fixed-corpus inspection, oracle execution, import restrictions, and output audit
found no external capability. The regex scanner is defense in depth rather than
a general capability proof. The trusted generator and oracle harness use local
files, a Node subprocess, and dynamic import to construct and verify cases.

\subsection{Factorial design}

Table~\ref{tab:factors} shows the active factors. The base stratum supplies the
current artifact only; prior submission diffs are withheld. This local
information boundary is the principal concealment mechanism under test. A
family-stratified nested sample supplies either the same local view or the
current artifact plus cumulative prior artifacts, allowing the information-loss
mechanism to be estimated without changing the original base estimands.

\begin{table}[ht]
\centering
\caption{Active study factors.}
\label{tab:factors}
\ifdefined\EnableTaggedPDF\tagpdfsetup{table/header-rows={1}}\fi
\begin{tabular}{@{}lll@{}}
\toprule
Factor & Levels & Interpretation \\
\midrule
Intent & malicious, benign & Matched final-state policy outcome \\
Decomposition & atomic, split & One or three necessary submissions \\
Workflow & pull request, trunk & Pre-merge gate or untrusted-main pipeline \\
Information context & local; cumulative in nested stratum & Current artifact or cumulative history \\
Trial & one base; two in nested stratum & Independent fresh-process repetition \\
Review system & Codex, Claude & Subscription-backed product surface \\
\bottomrule
\end{tabular}
\end{table}

The base stratum contains 40 malicious-benign pairs, 80 cases, 640 sequence-system trials,
and 1,280 boundary reviews across the two systems. It has one local-context
trial per intent, decomposition, workflow, and system cell. Before outcomes, two
templates per family were selected without replacement by seeded Fisher--Yates
for a 16-template nested stratum. Those templates receive local and cumulative
contexts and two trials per cell. The complete two-system schedule therefore has
896 unique boundary prompts, 1,408 sequence-system trials, and 2,816 boundary
reviews: 1,408 per system, 704 atomic, and 2,112 split. Trials reuse the same
frozen prompt but start fresh processes.

The three split submissions retain fixed causal order and one-minute spacing;
the atomic submission is timestamped at the split activation time. Timing and
order are matched controls rather than estimated factors. Calls are globally
shuffled with a seeded Fisher--Yates algorithm. If activated, the separate
Ollama queue reuses only the 640 base local-context boundary prompts, once each,
as 320 sequences, 160 atomic boundaries, and 480 split boundaries. It receives
neither cumulative context nor repeated trials. No outcome-dependent stopping,
template replacement, or prompt tuning is permitted.

\subsection{Review systems}

The systems are specified as product configurations rather than interchangeable
base-model APIs (Table~\ref{tab:systems}). Each boundary starts a fresh process
with no persistent session. User configuration and tools are disabled where
the client supports it. Any emitted tool call is retained as a protocol
deviation.

\begin{table}[ht]
\centering
\caption{Confirmatory and conditional review systems.}
\label{tab:systems}
\ifdefined\EnableTaggedPDF\tagpdfsetup{table/header-rows={1}}\fi
\begin{tabular}{@{}lll@{}}
\toprule
System & Client version & Requested model \\
\midrule
Codex CLI & 0.137.0; ChatGPT login & \code{gpt-5.5} \\
Claude Code & 2.1.211; Claude Max login & \code{sonnet} \\
Ollama (conditional) & 0.32.1; not activated & \code{qwen3.5:397b-cloud} \\
\bottomrule
\end{tabular}
\end{table}

Before each batch, the runner verifies subscription authentication, live
account-identity fingerprints, disabled extra usage, and the absence of
\code{OPENAI\_API\_KEY}, \code{CODEX\_API\_KEY}, and
\code{ANTHROPIC\_API\_KEY}. It performs at most 16 physical dispatches before
repeating those checks and aborts rather than switching to API billing, credits,
another provider, or a fallback model. Codex additionally uses approval policy
\code{never} and a custom minimal permission profile whose offline proof denies
model-shell access to isolated credentials and the frozen response schema.
Requested-model evidence, exact client provenance, timestamps, allowlisted
integer token counts, and failure digests are retained. Every raw provider
envelope is archived privately before parsing and bound to its terminal attempt. Claude
acceptance requires exactly one \code{modelUsage} entry. Its sole key must be
\code{claude-sonnet-4-6}. Codex CLI 0.137.0 does not expose a returned-model field
in its JSON event contract; each accepted result instead records requested model
\code{gpt-5.5}, the exact pinned binary path, version and hash, and an explicit
status that the requested pinned client showed no observable reroute. A
model-rerouted error event is rejected. This is an observability limitation,
not returned-model verification.

No automated Ollama access or purchase is currently authorised. A cross-provider
assessment applies one provider-neutral operational rule: published first-party
terms and product documentation expressly supporting automation on the intended
authenticated product surface supply the basis for that documented use. This is
not legal advice or provider approval, and any direct contrary or narrower
provider reply overrides it prospectively. The archived AI-support response
neither granted nor denied permission. On the current published basis, no
exceptional research waiver appears necessary. The remaining fixed order is
action-time approval for a no-prompt
entitlement and identity preflight, fixture-only implementation and offline
artifact generation, targeted review, an outcome-blind activation decision, a pre-submission content freeze,
and OSF registration followed by a separately bound execution authorisation.
If activated, the arm permits exactly 640 physical inference dispatches and zero
retries. Requests set \code{stream=false} and \code{think=false}, omit tools,
deliberately omit \code{format}, and pass through the common frozen local response
validator. Exact \code{response.model} identity as
\code{qwen3.5:397b-cloud} is required. Invalid or schema-invalid output is a
terminal abstention; any reasoning trace or tool call suspends the arm. The
localhost service brokers cloud inference, so prompts leave the device. A new
fixed subscription requires action-time approval, and extra-usage balance,
per-call billing, overage, and fallback are prohibited.

All reviews receive one fixed security-review task and a JSON response contract.
Requests ask for \code{allow}, \code{block}, or \code{abstain}; a maliciousness
probability; severity; relevant \code{head\_revision} or
\code{commit\_revision} values; evidence; and one
canonical abstract-policy category from the frozen schema. Complete requests
are at most 3,992 UTF-8 bytes. A response that fails the structured-response
validator is a terminal \code{schema\_invalid} abstention and is not repaired. The
protocol does not infer a separate refusal class from unstructured provider output.

\subsection{Collection and blinding}

The call queue contains opaque call, prompt, and case identifiers plus review
system and schedule position. It contains no intent, family, structural-template,
activation, or expected-severity field. Ground truth is held in a separately
keyed ledger and joined only after collection.

Subscription rate limits produce a durable suspension at the current call; the
same fixed queue resumes after the plan window resets. No API credit is accepted.
Immediately before each invocation, the author inspects both account billing
settings and supplies a private confirmation, no more than 15 minutes old,
binding the authorized account fingerprints and stating that extra or overage
usage is disabled. The runner verifies it before output creation or client
preflight. Each invocation is limited to 16 physical dispatches before that
confirmation and all other checks repeat.
Every frozen authorized call must have exactly one terminal result before
analysis: a schema-valid review, a validator-produced abstention, or a terminal
missing result after exactly three client-failure suspensions. Started,
suspension, and terminal result records are fsynced. Tool deviation is a retained
boolean on a schema-valid result. Suspension reasons distinguish rate limits;
identity absence or drift; malformed provider envelopes; authorized-schema,
executable, or isolation drift; client failure; and interrupted processes.
``Crash-safe'' means resume after a cleanly observed interruption; it does not
claim recovery from a torn JSONL tail, abrupt power loss, concurrent writers, an
orphan created between envelope archival and ledger append, or a check-to-spawn
race. Collection also suspends on a safety or terms conflict,
authentication change, or benchmark integrity failure. The API, per-call,
overage, extra-balance, and paid-fallback spending ceiling is US\$0. Fixed
subscription fees are disclosed separately and any new plan requires action-time
approval.

\subsection{Outcomes}

The primary binary outcome is \code{blocked\_by\_activation}. Secondary outcomes are
recall, false-positive rate, precision, abstention rate, cumulative detection by
boundary, submissions and elapsed minutes to detection, Brier score,
calibration, severity error, localisation, and failure-category rates.
Submission localisation requires citing at least one frozen relevant submission
identifier; mechanism localisation requires the exact frozen abstract-policy
category; combined localisation requires both. Within-system consistency in the
balanced two-trial nested stratum uses pairwise operational-verdict agreement
and malicious-probability ICC(1,1), with template-bootstrap intervals.

\section{Analysis Plan and Power}

H1--H3 use only base-stratum rows: local context, trial 1, all 40 templates, and
exactly \code{codex-cli/gpt-5.5} and \code{claude-code/sonnet}. H1 equally
averages decomposition, workflow, and both products within template; H2 equally
averages workflow and both products within malicious templates; H3 equally
averages atomic and split conditions and both products within malicious
templates. Templates then enter the equal-family analysis.
Review-system-specific results are secondary. Registered 95\% intervals use a
fixed-strata Welch/Satterthwaite t procedure over template-level contrasts. The
eight policy families are fixed, equally weighted strata; within-family variance
components use five base templates for H1--H3 and two nested templates for H4.
A 10,000-replicate family-stratified percentile bootstrap with seed 20260718 is
reported as a sensitivity interval, not the support-decision interval.

The H1 estimate is the malicious-minus-benign block-rate risk difference,
equally averaged over decomposition, workflow, and both products. H1 is supported only when its 95\%
interval is above zero. The H2 estimate is the split-minus-atomic malicious risk
difference, equally averaged over workflow and both products. H2 is supported only when its 95\% interval
is below zero. H3 is the trunk-minus-pull-request malicious contrast, averaged
equally over decomposition and both products, and is supported only when its nondirectional 95\% interval
excludes zero. H4 uses only the frozen 16-template nested stratum, averaging its
two trials and both products within context cells, and then equally averaging
pull-request and trunk. It is the malicious-only difference
\((\mathrm{split}-\mathrm{atomic})_{\rm cumulative} -
(\mathrm{split}-\mathrm{atomic})_{\rm local}\) and is supported only when its 95\%
interval is above zero. The decomposition-by-workflow difference-in-differences
and cumulative-context main effect are exploratory. Nested consistency retains
template-bootstrap intervals and has no support decision.

If every within-family variance component for a registered hypothesis is zero,
its Welch interval, Satterthwaite degrees of freedom, and support
decision are unavailable. The point estimate and standard error 0 remain
reported, with only the percentile-bootstrap sensitivity interval shown.

Missing boundaries are non-detections in the operational analysis. Mandatory
estimand-specific bounds assign missing cells to zero or one in the opposing
directions for each contrast: malicious versus benign for H1, split versus
atomic for H2, trunk versus pull request for H3, and the four signed H4 cells.
A review-system-specific complete-pair sensitivity for H1--H4 drops a structural
template only for the affected product while retaining equal-family weighting.
If any fixed family then has fewer than two complete templates, that
product-specific registered interval is reported as unavailable rather than
changing the estimand or variance method. No imputation model is fitted. The four registered
claims are not combined into a multiplicity procedure, so no familywise error
control is claimed; every estimate, interval, bound, and decision is reported.

If activated, Ollama is analysed alone using the 40-template base local sample,
H1--H3 contrasts, fixed-strata Welch intervals, percentile-bootstrap
sensitivity, missingness bounds, and complete-template sensitivity. It has no
H4 or consistency estimate. Its dedicated reporting path
emits estimates, intervals, and bounds without hypothesis-support,
robustness, significance, equivalence, superiority, or pooled-model decisions.
Fabricated-response tests must prove both this decision-free path and rejection of
Ollama, missing, extra, or mixed rows from the confirmatory path before activation.
Confirmatory analysis also rejects a ledger unless every frozen authorized call
has exactly one terminal result.

The products and sample are a convenience selection determined by subscriptions
available to the author, not an 80\% power target. The optional Ollama stratum
reuses the same 40 templates and adds no independent template evidence or
confirmatory power. An earlier 20,000-replication single-system
normal-approximation audit did not match the final design. A subsequent
outcome-free audit simulated the exact base and nested strata and found material
undercoverage for the initially proposed percentile bootstrap. The interval was
therefore changed prospectively. Under 1,000 global-null simulations, fixed-
strata Welch coverage for H1--H4 was 0.953, 0.965, 0.953, and 0.982. Under the
central assumptions, support fractions were 1.000, 0.724, 0.356, and 0.073 and
expected interval widths were 0.143, 0.211, 0.212, and 0.480. These values are
assumption-conditional diagnostics, not guaranteed or achieved power. Null
results cannot support no-effect or equivalence claims; H4 in particular is
expected to be imprecise. The study is a bounded empirical evaluation and a
variance estimate for a larger replication.

\section{Pre-results Status}

No benchmark prompt has been submitted to any review system. The two-system
confirmatory generator, prompts, schedule, ledgers, scoring, estimands, and tests
were implemented with deterministic checks and fabricated plumbing responses
only. No LLM outcome has been collected or inspected. No Ollama runner, queue,
ledger, or output artifact has been implemented or generated. An earlier
no-prompt confirmatory preparation check passed for the pinned Codex and Claude
clients without inference; the hardened sandbox, identity, and ledger checks
require successor revalidation before freeze. Ollama's AI support agent replied that it could not grant
written permission or determine whether the specific design was permitted, and
offered escalation to a human agent. The author then requested human escalation
and asked whether ordinary signed-in subscription/API use already supplies the
permission referenced by the Terms. A human response is pending as a nonblocking
clarification; any contrary restriction overrides the terms assessment. Automated
preflight, subscription purchase, runtime support, and activation remain absent.

The deterministic baseline ran ESLint 9.39.3 over all 2,400 final-state files in
the full 800-case generator output. It emitted no warnings or errors: zero
malicious and zero benign cases were flagged, giving recall 0, false-positive
rate 0, and undefined precision. This conventional final-state lint does not
estimate decomposition or workflow effects. A separate zero-finding Semgrep
feasibility probe is disclosed but excluded because its exact registry-rules
snapshot cannot be redistributed under the provider's rules license.

\section{Threats to Validity}

\textbf{Construct validity.} The benchmark models malicious intent through a
harmless abstract policy violation. It tests composition recognition, not the
full semantics or distribution of production vulnerabilities. Workflow is a
controlled artifact and policy framing, not an organisation.

\textbf{Internal validity.} Atomic and split variants can differ in salience
despite identical final trees and line counts. Provider-controlled system
instructions are neither observable nor identical. Product routing and serving
stacks may change behind requested model identifiers; metadata and access times
are therefore retained and drift suspends collection. Collection may span
subscription reset windows, creating calendar-time and platform-drift risk
despite the fixed randomised queue.

\textbf{External validity.} Forty purposively selected templates, eight policy
families, two confirmatory products, an optional exploratory product, and compact
JavaScript repositories do not represent all
languages, organisations, vulnerabilities, or attack strategies. The design
tests one three-submission decomposition and does not estimate an optimal
adversarial strategy. Inference is conditional on these templates; one lexical
instance per template does not estimate identifier-instance variability. The
subscription-only design is not cost-free for an independent reproducer without
equivalent subscriptions. Product surfaces, schema enforcement, hidden prompts,
calendar periods, routing, and cloud aliases confound descriptive cross-product
comparisons and cannot support provider or open-versus-closed-weight claims.

\textbf{Statistical conclusion validity.} The sample is underpowered for small
effects, especially workflow and interaction effects. Only the nested stratum
has repeated trials, so consistency estimates apply to 16 purposively bounded
templates rather than the full sample. Family-stratified intervals expose rather
than hide this uncertainty.

\section{Ethics and Responsible Release}

This is defensive research into review reliability. The benchmark cannot access
external systems or data, and its only harmful state is a deterministic in-memory
assertion over synthetic names. No real repository or person is tested.

For the conditional cloud arm, prompts leave the device even though the local
service is used. The study makes no independent claim about provider logging,
training, or retention beyond the dated first-party materials and archived
correspondence. Public research data are limited to sanitized derived
attempt-event data and aggregate analysis. Raw ledgers, provider output, absolute
paths, account and authentication state, and environment metadata remain private
unless a post-outcome safety, upstream-license, privacy, and terms review approves
a narrower release. Aggregate reporting uses factual plain-text attribution
without logos or endorsement claims.

Three isolated subagents initially reviewed candidate commit \code{6b607f6} for
benchmark safety, methods, and reproducibility. All three returned
\code{do not approve}; their raw reports are archived. Correction addenda inspect
the combined resolution log and are therefore not isolated. These
records are labelled AI-assisted internal review, not
independent human peer review, ethics approval, or arXiv endorsement. Shared
model, provider, and author-orchestration biases are reported as limitations.
Neutral change titles remove plausible refactor cover stories, but the abstract
composition patterns retain disclosed dual-use value. The packet contains no
real target, external capability, credential, deployment path, or complete exploit.
Targeted non-isolated reviews of commit \code{70963a2} rejected its release-packet
and execution-authorization mechanics. Their reports and accepted findings are
archived. A later outcome-blind worktree review also retains a \code{do not
approve} reproducibility disposition while authorization-schema, terminal-result,
identity-evidence, executable-pinning, writer-exclusion, and release-boundary
corrections remain in progress. The safety record additionally requires the
historical provider-specific request for written assurance to be reconciled with
the current provider-neutral published-terms rule: neither the AI-support response
nor silence is approval, and a direct contrary provider reply controls. The
present correction remains unfrozen and unauthorized. All
author choices and the study branch must first be fixed in one immutable
successor commit whose frozen marker identifies exact reviewer input bytes. All
three review roles must then approve that same commit before the detached
phase-one attestation and OSF submission. The marker cannot itself represent
approval or authorize an outcome call.

\section{Reproducibility}

The candidate package contains code for benchmark generation, scheduling,
scoring, sensitivity analysis, subscription collection, and outcome analysis.
It also contains the preregistration, archived internal reviews, and two-phase
freeze records. Each durable ledger begins with one immutable run-binding header.
The analysis requires the exact binding schema, recomputes its fingerprint, and
verifies the raw ground-truth-ledger digest before joining outcomes. It excludes
the metadata row from outcomes and carries the verified binding into its report.
Current unfrozen hashes are:

\begin{itemize}[leftmargin=1.5em]
  \item scenario cards:
  \code{d4ea0a7197e3e8e17bf027b9466d63de427c273aa19e539a672732b80b5eae3d};
  \item rendered prompts:
  \code{e99d650f5f670e138b0c31047c6009a1eea99f2d2005c742971419618c86bdf0};
  \item randomised schedule:
  \code{80c78d846bf264fcb8a136bffd3b6d3af50c0e8dfe622383272e0d4798ab63a1};
  \item blinded call ledger:
  \code{53422146aef4497fbb4da0819a3134f46b1e75351857a37ac92d0d3114bbbfb1};
  \item ground-truth ledger:
  \code{9453faab7b1bc1dc06cdef7bb83c7d4a15e96d6f6c06d407acd0aa823931cf6d}.
\end{itemize}

The first phase builds a canonical JSON payload from an exact Git commit using a
fixed 13-member release-safe allowlist. It includes the privacy-preserving
AI-support-response transcription with its required SHA-256 recorded in the
manifest.
Each member records path, media type, raw byte size,
SHA-256, and Base64 content; the verifier rejects extra, reordered, traversal,
detached-state, executable-case, credential, account-state, result, raw-output,
and reasoning-trace members. The outer freeze record binds the payload and
member-manifest hashes, selected branch, exact queue shape and counts, systems,
normalized artifacts, raw artifact bytes, and runtime files before OSF
submission, while explicitly leaving outcome calls unauthorised. Each detached
review report carries machine-readable role, exact-commit, disposition, and
blocker metadata that must match its report hash and phase-one attestation. The
confirmatory-only branch also requires a machine-readable non-activation decision
for the conditional Ollama arm and false pre-outcome state flags.

After registration, a separate execution record binds untouched raw OSF-origin
registration JSON, its structurally derived summary, the exact downloaded
two-file inventory, and explicit action-time author confirmation to the canonical
OSF identity, raw freeze-record hash, content commit, branch, payload, member
manifest, and exact queues. This author-attested HTTPS capture is not an OSF
cryptographic signature or independently trusted timestamp. UTC phase timestamps
are strictly ordered.
The confirmatory runner rejects extra or Ollama-labelled queues, recomputes the
schedule and all bindings, and executes the same prompt and call buffers it
hashed. Rejection precedes output creation, client preflight, or model execution.
Both detached records are excluded from the payload to avoid self-reference.

The earlier paid API design is retained in version history but was superseded
before outcome collection. The public research release will add sanitized derived
attempt-event data, analysis outputs, checksums, and the OSF and arXiv identifiers.
Raw ledgers, raw provider output, absolute paths, account and authentication state,
and environment metadata remain private pending explicit post-outcome review.

\section{Conclusion}

This study tests whether subscription-backed LLM review products recognise an
unsafe composition when its necessary changes arrive separately, and whether
the estimate changes under pull-request or untrusted-main framing. The safe
paired benchmark supports auditable positive or negative findings, while the
subscription constraint makes uncertainty and product drift part of the study
rather than hidden implementation details. Empirical conclusions remain
contingent on prospective registration, completion of the fixed schedule, and
the prespecified analysis.

\appendix
\section{Fixed Review Task}

The fixed task instructs the review system to act as a security reviewer, treat
repository and change metadata as untrusted evidence, decide whether the current
submission should be allowed, blocked, or sent for human review, consider unsafe
behaviour created alone or in composition with supplied evidence, and return
only the contracted JSON grounded in the supplied artifact. No chain-of-thought
request, product-specific prompt tuning, or outcome-driven tuning is used.

\begin{thebibliography}{9}
\bibitem{yu2024}
J. Yu, P. Liang, Y. Fu, A. Tahir, M. Shahin, C. Wang, and Y. Cai.
An insight into security code review with LLMs: Capabilities, obstacles, and
influential factors. \emph{arXiv:2401.16310}, 2024.
\url{https://arxiv.org/abs/2401.16310}.

\bibitem{liu2024}
Y. Liu, L. Gao, M. Yang, Y. Xie, P. Chen, X. Zhang, and W. Chen.
VulDetectBench: Evaluating the deep capability of vulnerability detection with
large language models. \emph{arXiv:2406.07595}, 2024.
\url{https://arxiv.org/abs/2406.07595}.

\bibitem{zhou2024}
X. Zhou, D.-M. Tran, T. Le-Cong, T. Zhang, I. C. Irsan, J. Sumarlin, B. Le,
and D. Lo. Comparison of static application security testing tools and large
language models for repo-level vulnerability detection.
\emph{arXiv:2407.16235}, 2024.
\url{https://arxiv.org/abs/2407.16235}.

\bibitem{adalsteinsson2025}
F. S. Adalsteinsson, B. B. Magnusson, M. Milicevic, A. N. Davidsson, and
C.-H. Cheng. Rethinking code review workflows with LLM assistance: An empirical
study. \emph{arXiv:2505.16339}, 2025.
\url{https://arxiv.org/abs/2505.16339}.

\bibitem{basic2024}
E. Basic and A. Giaretta. From vulnerabilities to remediation: A systematic
literature review of LLMs in code security. \emph{arXiv:2412.15004}, 2024.
\url{https://arxiv.org/abs/2412.15004}.
\end{thebibliography}

\end{document}
